\documentclass[12pt]{article}
\usepackage[spanish]{babel}
\usepackage[utf8]{inputenc}
\usepackage{enumitem}
\usepackage[a4paper, margin=2.5cm]{geometry}

\title{Cuestionario: Servicios de Red}
\date{}

\begin{document}

\maketitle

\section*{Instrucciones}
Seleccione la opción correcta en cada pregunta.

%------------------------------------------------
\section*{DHCP}

\begin{enumerate}[label=\textbf{DHCP-\arabic*.}]

\item ¿Cuál es la función principal de DHCP?
\begin{enumerate}[label=\alph*)]
\item Transferir archivos
\item Asignar direcciones IP automáticamente
\item Traducir dominios
\item Enviar correos
\end{enumerate}

\item ¿Qué significa DHCP?
\begin{enumerate}[label=\alph*)]
\item Dynamic Host Configuration Protocol
\item Data Host Control Protocol
\item Domain Host Configuration Program
\item Dynamic Hardware Communication Protocol
\end{enumerate}

\item ¿Qué información puede proporcionar un servidor DHCP?
\begin{enumerate}[label=\alph*)]
\item Dirección IP
\item Puerta de enlace
\item Servidor DNS
\item Todas las anteriores
\end{enumerate}

\item ¿En qué puerto trabaja DHCP principalmente?
\begin{enumerate}[label=\alph*)]
\item 22
\item 67/68
\item 80
\item 53
\end{enumerate}

\item ¿Qué problema resuelve DHCP en redes grandes?
\begin{enumerate}[label=\alph*)]
\item Configuración manual de IP
\item Transferencia lenta
\item Seguridad del correo
\item Compresión de datos
\end{enumerate}

\end{enumerate}

%------------------------------------------------
\section*{DNS}

\begin{enumerate}[label=\textbf{DNS-\arabic*.}]

\item ¿Cuál es la función principal del DNS?
\begin{enumerate}[label=\alph*)]
\item Convertir nombres de dominio en direcciones IP
\item Transferir archivos
\item Administrar usuarios
\item Compartir discos
\end{enumerate}

\item ¿Qué puerto usa DNS generalmente?
\begin{enumerate}[label=\alph*)]
\item 53
\item 22
\item 25
\item 443
\end{enumerate}

\item ¿Cuál es un ejemplo de servidor DNS público?
\begin{enumerate}[label=\alph*)]
\item 8.8.8.8
\item 192.168.1.1
\item 127.0.0.1
\item 10.0.0.1
\end{enumerate}

\item ¿Qué tipo de registro DNS asocia un dominio con una IP?
\begin{enumerate}[label=\alph*)]
\item MX
\item A
\item TXT
\item CNAME
\end{enumerate}

\item ¿Qué sucede si falla DNS?
\begin{enumerate}[label=\alph*)]
\item No se pueden resolver dominios
\item Se pierden archivos
\item Se bloquea la red
\item Se apaga el router
\end{enumerate}

\end{enumerate}

%------------------------------------------------
\section*{SSH}

\begin{enumerate}[label=\textbf{SSH-\arabic*.}]

\item ¿Para qué se usa SSH?
\begin{enumerate}[label=\alph*)]
\item Administración remota segura
\item Enviar correos
\item Transferir páginas web
\item Resolver dominios
\end{enumerate}

\item ¿Qué puerto utiliza SSH por defecto?
\begin{enumerate}[label=\alph*)]
\item 21
\item 22
\item 23
\item 80
\end{enumerate}

\item ¿Qué característica distingue a SSH de Telnet?
\begin{enumerate}[label=\alph*)]
\item Compresión
\item Cifrado
\item Uso de UDP
\item Multicast
\end{enumerate}

\item ¿Qué herramienta común usa SSH en Linux?
\begin{enumerate}[label=\alph*)]
\item ssh
\item ftp
\item curl
\item dig
\end{enumerate}

\item ¿Qué permite SSH además del acceso remoto?
\begin{enumerate}[label=\alph*)]
\item Túneles de red
\item Compartir impresoras
\item DHCP
\item NAT
\end{enumerate}

\end{enumerate}

%------------------------------------------------
\section*{FTP}

\begin{enumerate}[label=\textbf{FTP-\arabic*.}]

\item ¿Para qué sirve FTP?
\begin{enumerate}[label=\alph*)]
\item Transferencia de archivos
\item Correo electrónico
\item DNS
\item Autenticación
\end{enumerate}

\item ¿Qué puerto usa FTP?
\begin{enumerate}[label=\alph*)]
\item 21
\item 22
\item 25
\item 80
\end{enumerate}

\item ¿Qué problema tiene FTP tradicional?
\begin{enumerate}[label=\alph*)]
\item No tiene cifrado
\item Es muy lento
\item No usa TCP
\item No usa autenticación
\end{enumerate}

\item ¿Cuál es una alternativa segura a FTP?
\begin{enumerate}[label=\alph*)]
\item SFTP
\item SMTP
\item HTTP
\item SNMP
\end{enumerate}

\item ¿Qué permite FTP?
\begin{enumerate}[label=\alph*)]
\item Subir y descargar archivos
\item Resolver dominios
\item Asignar IP
\item Enviar correos
\end{enumerate}

\end{enumerate}

%------------------------------------------------
\section*{HTTP}

\begin{enumerate}[label=\textbf{HTTP-\arabic*.}]

\item ¿Qué protocolo usa la web?
\begin{enumerate}[label=\alph*)]
\item HTTP
\item FTP
\item DNS
\item SSH
\end{enumerate}

\item ¿Qué puerto usa HTTP?
\begin{enumerate}[label=\alph*)]
\item 80
\item 22
\item 25
\item 53
\end{enumerate}

\item ¿Qué versión segura existe?
\begin{enumerate}[label=\alph*)]
\item HTTPS
\item FTPS
\item SMTPS
\item DNSSEC
\end{enumerate}

\item HTTP funciona bajo el modelo:
\begin{enumerate}[label=\alph*)]
\item Cliente-servidor
\item Peer to peer
\item Multicast
\item Broadcast
\end{enumerate}

\item Un navegador web utiliza principalmente:
\begin{enumerate}[label=\alph*)]
\item HTTP
\item LDAP
\item NFS
\item DHCP
\end{enumerate}

\end{enumerate}

%------------------------------------------------
\section*{NFS}

\begin{enumerate}[label=\textbf{NFS-\arabic*.}]

\item ¿Qué permite NFS?
\begin{enumerate}[label=\alph*)]
\item Compartir sistemas de archivos en red
\item Enviar correos
\item Asignar IP
\item Resolver dominios
\end{enumerate}

\item NFS es común en sistemas:
\begin{enumerate}[label=\alph*)]
\item Linux/Unix
\item Windows exclusivamente
\item iOS
\item Android
\end{enumerate}

\item ¿Qué permite montar NFS?
\begin{enumerate}[label=\alph*)]
\item Directorios remotos
\item Correos
\item Usuarios
\item DNS
\end{enumerate}

\item NFS funciona sobre:
\begin{enumerate}[label=\alph*)]
\item RPC
\item SMTP
\item SNMP
\item POP3
\end{enumerate}

\item ¿Para qué se usa en servidores?
\begin{enumerate}[label=\alph*)]
\item Compartir almacenamiento
\item DNS
\item VPN
\item DHCP
\end{enumerate}

\end{enumerate}

%------------------------------------------------
\section*{LDAP}

\begin{enumerate}[label=\textbf{LDAP-\arabic*.}]

\item LDAP se usa para:
\begin{enumerate}[label=\alph*)]
\item Directorios de usuarios
\item Transferencia de archivos
\item DNS
\item HTTP
\end{enumerate}

\item LDAP significa:
\begin{enumerate}[label=\alph*)]
\item Lightweight Directory Access Protocol
\item Local Directory Access Program
\item Logical Data Access Protocol
\item Lightweight Data Application Protocol
\end{enumerate}

\item LDAP es común en:
\begin{enumerate}[label=\alph*)]
\item Active Directory
\item FTP
\item HTTP
\item DHCP
\end{enumerate}

\item LDAP organiza datos en forma de:
\begin{enumerate}[label=\alph*)]
\item Árbol jerárquico
\item Tabla
\item Grafo
\item Lista
\end{enumerate}

\item LDAP se usa para:
\begin{enumerate}[label=\alph*)]
\item Autenticación
\item Resolución DNS
\item Transferencia FTP
\item Streaming
\end{enumerate}

\end{enumerate}

%------------------------------------------------
\section*{SMTP}

\begin{enumerate}[label=\textbf{SMTP-\arabic*.}]

\item SMTP se usa para:
\begin{enumerate}[label=\alph*)]
\item Enviar correos
\item Descargar correos
\item DNS
\item FTP
\end{enumerate}

\item Puerto común de SMTP:
\begin{enumerate}[label=\alph*)]
\item 25
\item 22
\item 80
\item 53
\end{enumerate}

\item SMTP significa:
\begin{enumerate}[label=\alph*)]
\item Simple Mail Transfer Protocol
\item Secure Mail Transfer Program
\item System Mail Transmission Protocol
\item Simple Message Transport Program
\end{enumerate}

\item SMTP trabaja junto con:
\begin{enumerate}[label=\alph*)]
\item POP3 o IMAP
\item FTP
\item DNS
\item SNMP
\end{enumerate}

\item SMTP se usa principalmente en:
\begin{enumerate}[label=\alph*)]
\item Servidores de correo
\item Servidores web
\item Routers
\item Switches
\end{enumerate}

\end{enumerate}

%------------------------------------------------
\section*{Proxy}

\begin{enumerate}[label=\textbf{Proxy-\arabic*.}]

\item Un proxy actúa como:
\begin{enumerate}[label=\alph*)]
\item Intermediario entre cliente y servidor
\item DNS
\item Servidor DHCP
\item Switch
\end{enumerate}

\item Una función del proxy es:
\begin{enumerate}[label=\alph*)]
\item Filtrado de contenido
\item Asignar IP
\item Resolver dominios
\item Transferir correo
\end{enumerate}

\item Un proxy puede mejorar:
\begin{enumerate}[label=\alph*)]
\item Cache de contenido
\item DNS
\item DHCP
\item SMTP
\end{enumerate}

\item Un proxy ayuda a:
\begin{enumerate}[label=\alph*)]
\item Controlar acceso a internet
\item Crear discos
\item Enviar correos
\item DHCP
\end{enumerate}

\item Ejemplo de software proxy:
\begin{enumerate}[label=\alph*)]
\item Squid
\item Apache
\item Bind
\item Postfix
\end{enumerate}

\end{enumerate}

\end{document}
